\section{Explained-Variance Alignment Certificate}

This section records Supplementary Theorem S2.13, which rewrites the S2.12 residual penalty using explained-variance fractions.

Assume $U$ and $S$ are square-integrable, $S\ge0$, $0<\mathbb E[S]<\infty$, and both total variances are positive. Let
\[
m(Y)=\mathbb E[U\mid Y],
\qquad
a(Y)=\mathbb E[S\mid Y].
\]
Define
\[
A_U=
\frac{\operatorname{Var}(m(Y))}{\operatorname{Var}(U)},
\qquad
A_S=
\frac{\operatorname{Var}(a(Y))}{\operatorname{Var}(S)}.
\]

\begin{theorem}[Explained-variance alignment certificate]
The S2.12 bound can be written as
\[
\operatorname{Cov}(U,S)
\ge
\operatorname{Cov}(m(Y),a(Y))
-
\sqrt{\operatorname{Var}(U)\operatorname{Var}(S)(1-A_U)(1-A_S)}.
\]
If $A_UA_S>0$ and
\[
\rho_{ma}=\operatorname{Corr}(m(Y),a(Y)),
\]
then
\[
\operatorname{Cov}(U,S)
\ge
\sqrt{\operatorname{Var}(U)\operatorname{Var}(S)}
\left[
\rho_{ma}\sqrt{A_UA_S}
-
\sqrt{(1-A_U)(1-A_S)}
\right].
\]
Hence
\[
\rho_{ma}\sqrt{A_UA_S}
>
\sqrt{(1-A_U)(1-A_S)}
\]
is sufficient for positive total covariance.
\end{theorem}

\begin{proof}
The law of total variance gives
\[
\mathbb E[\operatorname{Var}(U\mid Y)]
=
\operatorname{Var}(U)-\operatorname{Var}(m(Y))
=
\operatorname{Var}(U)(1-A_U),
\]
and similarly
\[
\mathbb E[\operatorname{Var}(S\mid Y)]
=
\operatorname{Var}(S)(1-A_S).
\]
Substitution into the coarser S2.12 residual-variance bound gives the first inequality. If $A_UA_S>0$, then
\[
\operatorname{Cov}(m(Y),a(Y))
=
\rho_{ma}
\sqrt{A_UA_S\operatorname{Var}(U)\operatorname{Var}(S)},
\]
which yields the normalized form after factoring out the positive total standard-deviation product.
\end{proof}

\paragraph{Perfect conditional-mean alignment.}
If $\rho_{ma}=1$, the sufficient condition becomes
\[
A_UA_S>(1-A_U)(1-A_S),
\]
which simplifies exactly to
\[
A_U+A_S>1.
\]

\paragraph{Symmetric explained variance.}
If $A_U=A_S=A>0$, the condition becomes
\[
\rho_{ma}A>1-A.
\]
When $\rho_{ma}>-1$, this is algebraically equivalent to
\[
A>\frac{1}{1+\rho_{ma}}.
\]
Because $0<A\le1$, the strict certificate is feasible only when $\rho_{ma}>0$. For $\rho_{ma}=0$, it would require $A>1$; for $\rho_{ma}<0$, the threshold is even larger, and for $\rho_{ma}=-1$ the divided form is undefined. For $\rho_{ma}=1$, the threshold reduces to $A>1/2$.

\paragraph{Boundary.}
If $A_U=0$ or $A_S=0$, the normalized conditional-mean correlation may be undefined. The unnormalized theorem remains valid. The certificate is sufficient rather than necessary because S2.12 itself uses a worst-case residual anti-correlation penalty.
